% Logistic speciation --- plan \S3.3, "\S2".
% Sources: 09/03 entire, 09/04 entire, 09/02 (push of the past),
% 09/01:39-43 (the LSMean figure).
% Corrections C1 (09/04:7), C2 (09/04:34,35), C3 (09/04:176) applied here.
% Additions A1 (var:prop:yule-geom), A2 (var:cor:logistic-limit) and A3 (the
% beta_3 line) are new.

\section{Logistic speciation}
\label{var:sec:speciation}

\subsection{Variation, and what suppresses it}
\label{var:sec:variationmotivation}

The premise of the model is Fisher's: variation is in constant supply, and
what limits the number of distinct types is not the rate at which they arise
but the severity of the conditions under which they must persist
\cite{fisher1929genetical}. ``To maintain a stationary variance, fresh
mutations must be available in each generation'' is his statement of the
first half; the second half is that selection pressure, acting continuously,
holds a population near a current optimum and prevents most of that variation
from establishing. Nietzsche made the same observation from the breeder's
side, that species given superabundant nourishment ``tend in the most marked
way to develop variation''
\cite{nietzsche2013beyond}.
Read that way, an abundance of resources is less a cause of diversification
than a withdrawal of what had been suppressing it, and the \emph{speciation
rate} becomes a function of how much room is left.

That is a claim about the fossil record as much as about breeding. The
signature --- high early diversification after the opening of ecological
space, then a slowing --- is the one named by punctuated equilibria
\cite{eldredge1972punctuated}, and it is the one density-dependent
diversification models are built to produce \cite{morlon2014phylogenetic}.
Whether it is the typical pattern, and at what taxonomic level, is not
settled here.

\subsection{The push of the past}
\label{var:sec:potp}

There is an existing account of why diversification rates appear high at the
origin of a clade and decline towards the present, and the model below is a
response to it rather than a competitor. Mann and Budd, reviewing the
phenomenon introduced by Nee and co-workers \cite{Nee1994}, name it the push
of the past \cite{budd2018history}:

\begin{quotation}
\noindent
The POTPa emerges as a feature of diversification by the fact that all modern
clades (tautologically) survived until the present day. This singles them out
from the total population of clades that could be generated from any
particular pair of background birth and death rates: clades that happened by
chance to start off with higher net rates of diversification have
better-than-average chances of surviving until the present day. As Nee et
al.\ (1994a) put it, such clades ``got off to a flying start''; and they
accumulate species faster than one would expect. Once clades become
established, they are less vulnerable to random changes in the observed
diversification rate, and this value therefore tends to revert to the
background rate through time. Long-lived clades thus tend to show high
observed rates of diversification at their origin, which then decrease to
their long-term average as the present is approached.
\end{quotation}

\noindent
The push of the past is a survivorship effect: the rates never move, and the
apparent decline is an artefact of conditioning on survival to the present.
The model of this section supposes instead that the rate really does move,
declining because the ecological room that permitted the early
diversification is used up. The two accounts are not exclusive and nothing
here argues against the first. What the model offers is a second mechanism
with the same signature, and a distribution against which it might one day be
told apart from the first \cite{bak1991self}.

\subsection{The model}
\label{var:sec:speciationmodel}

Let $N(t)$ be the total abundance of individuals across all species, taken as
a real number on the grounds that it is large compared with the number of
species, and let it grow logistically,
\begin{equation}
\ddt{N}=\varrho N\lrb{1-\frac{N}{K}},
\qquad N(0)=N_0,\qquad 0<N_0<K,
\label{var:eq:logistic}
\end{equation}
with $K$ the carrying capacity of the environment and $\varrho$ the growth
rate. Its solution is
\begin{equation}
N(t)=\frac{K\ee^{\varrho t}}{\xi+\ee^{\varrho t}},
\qquad
\xi=\frac{K}{N_0}-1.
\label{var:eq:logisticsol}
\end{equation}

Let $S_t$ count species, with $S_0=1$. Each existing species gives rise to a
new one independently at the per-species rate $\beta(t)$, and no species
becomes extinct, so the transitions in a vanishing interval are
\begin{align}
\pr{S_{t+\Delta t}-S_t=1\mid S_t=n}&=\beta(t)\,n\,\Delta t+o(\Delta t),
\label{var:eq:transitions}\\
\pr{S_{t+\Delta t}-S_t=0\mid S_t=n}&=1-\beta(t)\,n\,\Delta t+o(\Delta t).
\notag
\end{align}
This is the Yule process of \cref{m:sec:yule} with its constant rate replaced
by a function of time; setting $\beta(t)=\beta$ recovers it exactly. The
substitution is not new in itself --- taking $\beta(t)=\ee^{-t}$ makes the
mean of $S_t$ follow Gompertz's 1825 law \cite{gompertz1825xxiv} --- and the
Yule process was itself introduced in 1925 to model exactly this quantity,
the number of species within a genus \cite{Yule1925}. What is being tried
here is a rate tied to the logistic environment:
\begin{equation}
\beta(t)=\sigma\bigl(K-N(t)\bigr),
\label{var:eq:rate1}
\end{equation}
which by \cref{var:eq:logisticsol} is
\begin{equation}
\beta(t)=\sigma K\,\frac{\xi}{\xi+\ee^{\varrho t}},
\label{var:eq:ratesolved}
\end{equation}
decreasing from $\sigma(K-N_0)$ to zero as the environment fills. The
coupling runs one way only: the environment sets the speciation clock, and
$S_t$ does not act back on $N(t)$. The same model appears as a worked example
in \cref{m:sec:logisticspeciation}, where its mean is recorded; what is added
here is its distribution.

\begin{remark}[What the rate law is, and is not, claiming]
\label{var:rem:scope}
\Cref{var:eq:rate1} treats resource availability as an \emph{index} of
selective pressure and not as a mechanism. It says that when there is room,
variants persist more often, and it says nothing about how. The model is
correspondingly agnostic about the direction of causation: it does not
distinguish logistic expansion driving diversification from logistic
expansion merely creating the space in which diversification happens.

Underneath that sits a question of timescale. If population expansion is
orders of magnitude faster than speciation --- which for most plausible
parameter values it is --- then the logistic phase is over before the species
count has moved appreciably, and \cref{var:eq:rate1} is better read as
setting the initial condition for a density-dependent diversification model
than as driving the rate through time \cite{morlon2014phylogenetic}. Nothing
below depends on resolving that, since the mathematics of
\cref{var:prop:yule-geom} holds for any locally integrable $\beta$, and the
question is which $\beta$ the biology warrants.
\end{remark}

\subsection{The marginal law}
\label{var:sec:marginal}

The first result is that the rate law enters the distribution of $S_t$ only
through its integral, and that the resulting law is geometric. Write
\begin{equation}
\Bint(t)=\int_0^t\beta(u)\,\dd u
\label{var:eq:Bdef}
\end{equation}
for the integrated speciation rate, and assume only that $\beta$ is locally
integrable and non-negative.

\begin{proposition}[Marginal law of an inhomogeneous Yule process]
\label{var:prop:yule-geom}
Let $S_t$ be a pure birth process on $\{1,2,\dots\}$ with $S_0=1$ and
time-dependent per-capita rate $\beta(t)$. Then $S_t$ is geometric on
$\{1,2,\dots\}$ with parameter $\ee^{-\Bint(t)}$:
\begin{equation}
\pr{S_t=n}=\ee^{-\Bint(t)}\bigl(1-\ee^{-\Bint(t)}\bigr)^{n-1},
\qquad n\ge1,
\label{var:eq:geomlaw}
\end{equation}
with
\[
\ex{S_t}=\ee^{\Bint(t)},
\qquad
\var{S_t}=\ee^{\Bint(t)}\bigl(\ee^{\Bint(t)}-1\bigr).
\]
\end{proposition}

The proof is the generating-function computation of
\cref{var:app:pgf}: the forward Kolmogorov equation still yields to the
method of characteristics, and the \pgf{} is that of a geometric law on
$\{1,2,\dots\}$ with success parameter $\ee^{-\Bint(t)}$. Three things are
worth recording here. The first is that the constant-rate Yule
law \cref{m:eq:yulelaw} is the case $\Bint(t)=\beta t$, so nothing is being
rebuilt: the substitution $\beta t\mapsto\Bint(t)$ is exact and not merely a
convenience. The second is that the argument nowhere uses the form of
$\beta$, only that its integral exists, so two rate laws with the same
integral to time $t$ give the same distribution at $t$ however differently
they arrived there. The third is that the variance falls out of the same
line, and had not been recorded before. The standard deviation is
$\ee^{\Bint}\sqrt{1-\ee^{-\Bint}}$, of the same order as the mean, so to fit
this model to a single lineage is to fit a mean to one draw from a law whose
coefficient of variation is close to one.

\subsection{The mean, directly}
\label{var:sec:mean}

With the proposition in hand $\ex{S_t}=\ee^{\Bint(t)}$ is a corollary, and it
is also \cref{m:eq:meanodeinhom} with $\mu\equiv0$. The direct derivation is
retained as a check on the integral that follows.

Over a vanishing interval the mean after is the mean before plus the expected
change, and by \cref{var:eq:transitions} the expected change is
$\beta(t)\ex{S_t}\Delta t+o(\Delta t)$. Dividing by $\Delta t$ and letting it
vanish,
\begin{equation}
\ddt{}\ex{S_t}=\beta(t)\ex{S_t},
\label{var:eq:meanode}
\end{equation}
which separates. With $\ex{S_0}=1$ the solution is $\ex{S_t}=\ee^{\Bint(t)}$,
and for the logistic rate law \cref{var:eq:ratesolved} the integral is
elementary:
\begin{equation}
\int\frac{\dd t}{\xi+\ee^{\varrho t}}
=\frac{1}{\xi}\lrb{t-\frac{1}{\varrho}\ln\lrb{\xi+\ee^{\varrho t}}},
\label{var:eq:intbeta}
\end{equation}
so that
\begin{equation}
\Bint(t)=\sigma K\lrb{t-\frac{1}{\varrho}\ln\lrb{\xi+\ee^{\varrho t}}
+\frac{1}{\varrho}\ln\lrb{\xi+1}},
\qquad
\ex{S_t}=\ee^{\sigma Kt}
\lrb{\frac{\xi+1}{\xi+\ee^{\varrho t}}}^{\sigma K/\varrho}.
\label{var:eq:Bt}
\end{equation}

\subsection{The limit law}
\label{var:sec:limit}

The question the model was built to answer is what happens at late time, when
$\beta$ has fallen to nothing and the species count has stopped moving. Since
$S_t$ is non-decreasing and its mean is bounded, it converges almost surely
to a finite limit; the content of the following is that the limit is proper
and its law is known.

\begin{corollary}[Limit law under the logistic rate]
\label{var:cor:logistic-limit}
Under \cref{var:eq:rate1}, $\Bint(t)$ converges as $t\to\infty$, to
\begin{equation}
\Bint(\infty)=\frac{\sigma K}{\varrho}\ln\lrb{\frac{K}{N_0}},
\label{var:eq:Binf}
\end{equation}
and $S_\infty$ is geometric on $\{1,2,\dots\}$ with
\begin{equation}
\ex{S_\infty}=\lrb{\frac{K}{N_0}}^{\sigma K/\varrho}.
\label{var:eq:limitmean}
\end{equation}
\end{corollary}

\begin{proof}
From \cref{var:eq:Bt},
\[
t-\frac{1}{\varrho}\ln\lrb{\xi+\ee^{\varrho t}}
=-\frac{1}{\varrho}\ln\lrb{1+\xi\ee^{-\varrho t}}
\;\longrightarrow\;0 ,
\]
so $\Bint(t)\to(\sigma K/\varrho)\ln(\xi+1)$, and $\xi+1=K/N_0$. The law is
\cref{var:eq:geomlaw} with $\Bint(t)$ replaced by its limit, which is
legitimate because $S_t$ increases to $S_\infty$ and the geometric
probabilities are continuous in the parameter.
\end{proof}

The dependence on the parameters is now explicit, and it is simpler than the
five constants $\sigma,K,\varrho,N_0,\xi$ suggest: only two dimensionless
groups survive. The first is $K/N_0$, the room the genus had at the outset,
so that a clade founded into a nearly empty environment integrates more
unused capacity than one founded into a full one and ends with more species.
The second is $\sigma K/\varrho$, the speciation capacity per unit of
population growth rate, which sets how much of that room is converted before
it closes; doubling the growth rate and the speciation coefficient together
changes nothing. That is the answer to the question posed but not settled
when the model was written --- how the asymptotic behaviour depends on the
parameters and their interplay --- and it is an answer about the whole
distribution rather than about the mean alone.

\begin{figure}[!htbp]
\centering
\includegraphics[width=0.92\textwidth]{fig02b_logspec_limit/fig02b}
\caption{What \cref{var:cor:logistic-limit} says, drawn from the closed form.
Panel~(a): the law of $S_\infty$ is geometric, so most of its mass sits at
small counts however large the mean, and its coefficient of variation is close
to one --- fitting this model to a single clade is fitting a mean to one draw
from that. Panel~(b): the expected final count
$\ex{S_\infty}=(K/N_0)^{\sigma K/\varrho}$ depends on the five constants
$\sigma,K,\varrho,N_0,\xi$ only through the two dimensionless groups on the
axes, so every point on a contour is a genus with the same expected diversity
reached from different parameters.}
\label{var:fig:logspec-limit}
\end{figure}

\begin{figure}[!htbp]
\centering
\includegraphics[width=0.72\textwidth]{fig02_logspec_mean/fig02}
\caption{Mean species count \cref{var:eq:Bt} for $N_0=1,2,5$, with $K=10$,
$\varrho=1$ and $\sigma=0.08$; dashed lines are the limits
$(K/N_0)^{\sigma K/\varrho}$. A smaller founding abundance leaves more
integrated unused capacity, so the limiting mean is higher. This is the
pure-birth model of \cref{var:eq:transitions}, not the birth--death numerical
variant recorded at \cref{m:fig:lsmean};
\cref{var:sec:withextinction} is what a constant extinction rate does to this
one.}
\label{var:fig:lsmean}
\end{figure}

\subsection{The other rate laws}
\label{var:sec:altrates}

The form \cref{var:eq:rate1} was one of five candidates written down when the
model was set up. The others were
\begin{align}
\beta_2(t)&=\sigma\bigl(K-N(t)\bigr)+\eta, &
\beta_3(t)&=\sigma\,\ddt{N}, \notag\\
\beta_4(t)&=\sigma\,\ddt{N}+\eta, &
\beta_5(t)&=\sigma\bigl(K_t-N(t)\bigr),
\label{var:eq:altrates}
\end{align}
with $\eta$ a constant floor and $K_t$ a capacity switching at random. The
motivation for $\beta_3$ and $\beta_4$ is that a mutation event needs a
replication event, so speciation might track the rate of population increase
rather than the room remaining; the floor $\eta$ in $\beta_2$ and $\beta_4$
is a baseline propensity to vary that survives the environment filling up.
\Cref{var:prop:yule-geom} settles two of the four immediately, because it
asks only for an integral.

For $\beta_3$ the integral is the total change in abundance:
\begin{equation}
\Bint(\infty)=\sigma\int_0^\infty N'(u)\,\dd u=\sigma\,(K-N_0),
\label{var:eq:beta3}
\end{equation}
so $S_\infty$ is again geometric, with mean $\ee^{\sigma(K-N_0)}$. The
comparison with \cref{var:eq:limitmean} is worth a moment. Under $\beta_1$
the limiting mean is a power of $K/N_0$ and under $\beta_3$ an exponential of
$K-N_0$, so the first is sensitive to the \emph{ratio} of capacity to
founding abundance and the second to their \emph{difference}; a model fitted
under one and read under the other would be badly wrong for large $K$.

For $\beta_2$ and $\beta_4$ the floor makes $\Bint(t)\sim\eta t$, which
diverges, and there is then no limit law: $S_t$ remains geometric at every
finite time with parameter $\ee^{-\Bint(t)}\to0$, and the species count grows
without bound, eventually at the constant-rate Yule rate $\eta$. That is a
statement about the model, not a defect of the analysis: a floor on the
speciation rate says diversification never stops, and a process that never
stops has no final species count for a distribution to be about.

$\beta_5$ is the one that remains genuinely open, and the proposition says
why. If $K_t$ switches at random then $\Bint(t)$ is itself a random variable,
and \cref{var:eq:geomlaw} holds conditionally on it. Unconditionally $S_t$ is
a \emph{mixture} of geometric laws over the law of $\Bint(t)$,
\begin{equation}
\pr{S_t=n}=\ex{\ee^{-\Bint(t)}\bigl(1-\ee^{-\Bint(t)}\bigr)^{n-1}},
\label{var:eq:mixture}
\end{equation}
so that the problem becomes finding the law of the integrated rate under the
switching regime. For a finite-state Markov switching $K_t$ that is a
Markov-modulated integral and is tractable; for anything else it is not
obviously so. Of five candidate rate laws, two are now closed, two are
explained non-results, and one is a well-posed problem.

The pattern across all five is the same one: nothing about a moving
speciation rate is difficult so long as its integral is deterministic, and
everything about it is difficult once the integral is not.

\subsection{What extinction does}
\label{var:sec:withextinction}

Everything above depends on the process being pure birth.
\Cref{var:eq:transitions} has no downward channel, and both the geometric
marginal of \cref{var:prop:yule-geom} and the limit law of
\cref{var:cor:logistic-limit} come of the fact that $\Bint$ converges and
nothing is ever taken away. It is worth seeing what the same origination law
does once species are allowed to go extinct, because the answer is not a
slower version of the same thing.

Let each species become extinct independently at the constant per-capita rate
$\mu$, so that $S_t$ is a linear birth--death process with a moving birth
rate and a fixed death rate. Its mean is the exponential of the net
integrated rate,
\begin{equation}
\ex{S_t}=\exp\lrb{\Bint(t)-\mu t},
\label{var:eq:meanwithdeath}
\end{equation}
which follows from $\dd\ex{S_t}/\dd t=(\beta(t)-\mu)\ex{S_t}$ and needs no
generating function. Two things follow at once. The mean turns over exactly
where the two rates cross, at the $t$ solving $\beta(t)=\mu$, whatever the
rate law; and since $\Bint(\infty)$ is finite under every law of
\cref{var:sec:altrates} while $\mu t$ is not, $\ex{S_t}\to0$. A constant
extinction rate set against an origination rate that runs out takes the clade
with it, and there is no limit law of any kind.

\Cref{var:fig:extinction} draws that under $\beta_3(t)=\sigma N'(t)$, with
one modification: the extinction of a genus's last species is suppressed, so
that the picture is about the shape of the rise and fall rather than about
the disappearance of the clade. The peak is unaffected, sitting at the
crossing as \cref{var:eq:meanwithdeath} requires, and the mean settles at one
species rather than at zero.

What the figure shows is the signature
\cref{var:sec:variationmotivation} began with, and the pure-birth model of
\cref{var:sec:speciationmodel} does not produce it. There the species count
rises and saturates, because the rate runs out and nothing erodes what it
built; \cref{var:cor:logistic-limit} is exactly that statement. Here it
rises, turns over and decays, which is the shape the fossil record is usually
described as having. The two accounts of \cref{var:sec:potp} are worth
holding against this. The push of the past produces a declining \emph{apparent}
rate by conditioning on survival while the true rates stay fixed;
\cref{var:eq:rate1} produces one by moving the rate and keeping every clade.
\Cref{var:fig:extinction} is a third thing again --- rates that move
\emph{and} a real loss --- and it is the only one of the three in which the
count itself goes down.

\begin{figure}[!htbp]
\centering
\includegraphics[width=0.84\textwidth]{fig02c_extinction/fig02c}
\caption{Logistic speciation with a constant extinction rate, at
$(K,N_0,\varrho,\sigma,\mu)=(10,5,0.8,0.8,0.4)$ from $S_0=1$: $160$ exact
realisations in pale blue, their mean over $400$ paths in bold. Origination
follows $\beta_3(t)=\sigma N'(t)$, drawn below against the constant $\mu$.
$N_0=K/2$ starts the abundance at the inflection of
\cref{var:eq:logistic}, so origination is at its largest at $t=0$ and decays
from there. The mean peaks at $10.0$ species at $t=3.24$, and the rates cross
at $t=3.29$: the two coincide, as \cref{var:eq:meanwithdeath} requires.
Extinction of a genus's last species is suppressed, which is why the paths
have a floor at one and the mean settles at $1.02$ rather than decaying to
zero. Set against the same origination law without extinction, where
\cref{var:eq:beta3} gives a limiting mean of
$\ee^{\sigma(K-N_0)}\approx55$ approached from below and never left, the cost
of a constant $\mu$ is the whole limit law.}
\label{var:fig:extinction}
\end{figure}

\paragraph{Open questions.}\label{var:rem:open-speciation}
\begin{enumerate}[leftmargin=1.6em,itemsep=0.2em,topsep=0.2em]
\item \textbf{The mixture at $\beta_5$.} What is the law of $\Bint(t)$ when
$K_t$ switches between finitely many values at exponential rates, and what
does the resulting mixture \cref{var:eq:mixture} do to the tail of $S_t$?
This is the question in this section most worth taking up, and it asks for a
Markov-modulated integral rather than for a new model.
\item \textbf{Identifiability.} Since only $\Bint(t)$ enters
\cref{var:eq:geomlaw}, no observation of $S_t$ at a single time can separate
$\sigma$, $K$ and $\varrho$ beyond the two groups of
\cref{var:eq:limitmean}. What a time series of $S_t$ at several times does
identify, and whether $\beta_1$ and $\beta_3$ can be told apart from one, is
not known.
\item \textbf{The timescale question of \cref{var:rem:scope}.} Whether
population expansion and speciation are separated by orders of magnitude in
any clade for which both are measurable, and therefore whether the model
describes a rate or an initial condition, is an empirical question this
chapter cannot settle.
\item \textbf{The marginal law under extinction.} \Cref{var:eq:meanwithdeath}
gives the mean but says nothing about the distribution, and
\cref{var:prop:yule-geom} does not survive a downward channel as it stands.
Whether the geometric marginal survives with an atom at zero --- so that
\cref{var:prop:yule-geom} is the $\mu=0$ case of something more general, and
the rate law still enters only through its integral --- is the question
\cref{var:fig:extinction} raises and does not answer.
\end{enumerate}

\FloatBarrier
